\section{Risks and Limitations}

The main technical risk is scope expansion. Phase 2 can easily become too large if it attempts multi-source radar, full MUSIC implementation, fixed-point quantization, and hardware deployment at the same time. To control this risk, Phase 2 will prioritize reproducible robustness experiments before optional extensions.

\begin{riskbox}{Scope Control Rule}
Phase 2 must first prove that the Phase 1 pipeline can be reproduced and stress-tested. FPGA, PYNQ, and Vitis HLS implementation should remain future phases unless the core Phase 2 results are already stable.
\end{riskbox}

Additional limitations include:

\begin{itemize}
    \item Synthetic data may not fully represent real radar measurements.
    \item Sector classification is less ambitious than continuous high-resolution DOA estimation.
    \item Rate-based spike encoding may introduce stochastic variability.
    \item LIF-inspired accumulation is a simplified temporal inference model, not a full biologically realistic SNN.
    \item MUSIC or multi-source extensions may require careful interpretation and should not be overclaimed.
\end{itemize}

These limitations are acceptable because the objective of Phase 2 is not to claim a final radar product. The objective is to build a stronger scientific bridge from Phase 1 toward later embedded and hardware-aware implementation.
